% GENERATED FILE. Edit canonical lessons, then run npm run generate:books.
% canonical-source-hash: fnv1a64:d3e60fa3ce56f785
% canonical-lessons: JA-C304-orei, JA-C304-ongaeshi, JA-C304-wana, JA-C304-nozoku, JA-C304-mazushii

\chapter{Repaying a Kindness}
\label{ch:ja-repaying-a-kindness}
\hlchaptermodality{304}

\begin{chapteropening}
\textbf{By the end of this chapter:} \emph{I can tell how a rescued bird repaid the old couple.}

Complete the last lesson of chapter 304: I can tell how a rescued bird repaid the old couple.
\end{chapteropening}

\section[\texorpdfstring{\ja{おれい} (orei)}{おれい (orei)}]{\ja{おれい} (orei) — thanks; a thank-you gift}

\label{lesson:JA-C304-orei}



\begin{quote}
Before the new one: say the Japanese for a small wayside shrine, then the Japanese for auspicious.
\end{quote}

\subsection*{You'll want to know: \ja{おれい}}
\textbf{\ja{おれい}} — \emph{orei} — \textquotedblleft{}thanks; a thank-you gift\textquotedblright{}.

\begin{quote}
\textbf{\ja{ははに} \ja{おれいを} \ja{いいました。}} — \emph{haha ni orei o iimashita} — \textquotedblleft{}I said thank you to my mother.\textquotedblright{}
\end{quote}

\begin{cousinweb}
\textbf{\ja{おれい}} (\emph{orei}) is a mixed word, native and Sino-Japanese: \emph{o-}, \textquotedblleft{}the polite prefix\textquotedblright{}, is native Japanese, and \emph{rei}, \textquotedblleft{}thanks\textquotedblright{}, is a Chinese character read with its on'yomi, a reading borrowed from Chinese (Middle Chinese \emph{lejX}) (JMdict, KANJIDIC2; Middle Chinese in Baxter's transcription).
\end{cousinweb}

\begin{grammarlens}[title={What you've built}]
One more word for telling what happened.
\end{grammarlens}

\subsection*{Your turn}
\begin{itemize}
\raggedright
  \item \emph{Say it:} \emph{orei}
  \item \emph{Say it:} \emph{orei}, once more
  \item \emph{Say it:} \emph{haha ni orei o iimashita} — I said thank you to my mother.
  \item \emph{Recall:} say the Japanese for a small wayside shrine, then the Japanese for auspicious, then say \emph{orei} again
\end{itemize}

\begin{tcolorbox}[breakable,colback=teal!4,colframe=teal!35!black,title={Before you move on}]
What does \ja{おれい} mean? (\textquotedblleft{}Thanks; a thank-you gift\textquotedblright{}.) Say it once more.
\end{tcolorbox}
\section[\texorpdfstring{\ja{おんがえし} (ongaeshi)}{おんがえし (ongaeshi)}]{\ja{おんがえし} (ongaeshi) — repaying a kindness}

\label{lesson:JA-C304-ongaeshi}



\begin{quote}
Before the new one: say the Japanese for auspicious, then the Japanese for thanks.
\end{quote}

\subsection*{You'll want to know: \ja{おんがえし}}
\textbf{\ja{おんがえし}} — \emph{ongaeshi} — \textquotedblleft{}repaying a kindness\textquotedblright{}.

\begin{quote}
\textbf{\ja{その} \ja{とりは} \ja{おじいさんに} \ja{おんがえしを} \ja{しました。}} — \emph{sono tori wa ojiisan ni ongaeshi o shimashita} — \textquotedblleft{}The bird repaid the old man's kindness.\textquotedblright{}
\end{quote}

\begin{cousinweb}
\textbf{\ja{おんがえし}} (\emph{ongaeshi}) is a mixed word, native and Sino-Japanese: \emph{on}, \textquotedblleft{}kindness\textquotedblright{}, is a Chinese character read with its on'yomi, a reading borrowed from Chinese (Middle Chinese \emph{'on}), and \emph{gaeshi}, \textquotedblleft{}return\textquotedblright{}, is native Japanese (JMdict, KANJIDIC2; Middle Chinese in Baxter's transcription).
\end{cousinweb}

\begin{grammarlens}[title={What you've built}]
One more word for telling what happened.
\end{grammarlens}

\subsection*{Your turn}
\begin{itemize}
\raggedright
  \item \emph{Say it:} \emph{ongaeshi}
  \item \emph{Say it:} \emph{ongaeshi}, once more
  \item \emph{Say it:} \emph{sono tori wa ojiisan ni ongaeshi o shimashita} — The bird repaid the old man's kindness.
  \item \emph{Recall:} say the Japanese for auspicious, then the Japanese for thanks, then say \emph{ongaeshi} again
\end{itemize}

\begin{tcolorbox}[breakable,colback=teal!4,colframe=teal!35!black,title={Before you move on}]
What does \ja{おんがえし} mean? (\textquotedblleft{}Repaying a kindness\textquotedblright{}.) Say it once more.
\end{tcolorbox}
\section[\texorpdfstring{\ja{わな} (wana)}{わな (wana)}]{\ja{わな} (wana) — a trap; a snare}

\label{lesson:JA-C304-wana}



\begin{quote}
Before the new one: say the Japanese for thanks, then the Japanese for repaying a kindness.
\end{quote}

\subsection*{You'll want to know: \ja{わな}}
\textbf{\ja{わな}} — \emph{wana} — \textquotedblleft{}a trap; a snare\textquotedblright{}.

\begin{quote}
\textbf{\ja{おじいさんは} \ja{わなの} \ja{とりを} \ja{たすけました。}} — \emph{ojiisan wa wana no tori o tasukemashita} — \textquotedblleft{}The old man rescued the bird from the trap.\textquotedblright{}
\end{quote}

\begin{cousinweb}
\textbf{\ja{わな}} (\emph{wana}) is a native Japanese word (wago), inherited rather than borrowed from Chinese; its kanji was attached later as a kun'yomi (JMdict).
\end{cousinweb}

\begin{grammarlens}[title={What you've built}]
One more word for telling what happened.
\end{grammarlens}

\subsection*{Your turn}
\begin{itemize}
\raggedright
  \item \emph{Say it:} \emph{wana}
  \item \emph{Say it:} \emph{wana}, once more
  \item \emph{Say it:} \emph{ojiisan wa wana no tori o tasukemashita} — The old man rescued the bird from the trap.
  \item \emph{Recall:} say the Japanese for thanks, then the Japanese for repaying a kindness, then say \emph{wana} again
\end{itemize}

\begin{tcolorbox}[breakable,colback=teal!4,colframe=teal!35!black,title={Before you move on}]
What does \ja{わな} mean? (\textquotedblleft{}A trap; a snare\textquotedblright{}.) Say it once more.
\end{tcolorbox}
\section[\texorpdfstring{\ja{のぞく} (nozoku)}{のぞく (nozoku)}]{\ja{のぞく} (nozoku) — to peep; to peek}

\label{lesson:JA-C304-nozoku}



\begin{quote}
Before the new one: say the Japanese for repaying a kindness, then the Japanese for a trap.
\end{quote}

\subsection*{You'll want to know: \ja{のぞく}}
\textbf{\ja{のぞく}} — \emph{nozoku} — \textquotedblleft{}to peep; to peek\textquotedblright{}.

\begin{quote}
\textbf{\ja{おばあさんは} \ja{へやの} \ja{なかを} \ja{のぞきました。}} — \emph{obāsan wa heya no naka o nozokimashita} — \textquotedblleft{}The old woman peeped into the room.\textquotedblright{}
\end{quote}

Polite \textbf{\ja{のぞきます}}, \emph{nozokimasu}; a godan verb. Past (polite): \textbf{\ja{おとこが} \ja{へやを} \ja{のぞきました。}}, \emph{otoko ga heya o nozokimashita}.

\begin{cousinweb}
\textbf{\ja{のぞく}} (\emph{nozoku}) is a native Japanese word (wago), inherited rather than borrowed from Chinese; its kanji was attached later as a kun'yomi (JMdict).
\end{cousinweb}

\begin{grammarlens}[title={What you've built}]
One more word for telling what happened.
\end{grammarlens}

\subsection*{Your turn}
\begin{itemize}
\raggedright
  \item \emph{Say it:} \emph{nozoku}
  \item \emph{Say it:} \emph{nozoku}, once more
  \item \emph{Say it:} \emph{obāsan wa heya no naka o nozokimashita} — The old woman peeped into the room.
  \item \emph{Recall:} say the Japanese for repaying a kindness, then the Japanese for a trap, then say \emph{nozoku} again
\end{itemize}

\begin{tcolorbox}[breakable,colback=teal!4,colframe=teal!35!black,title={Before you move on}]
What does \ja{のぞく} mean? (\textquotedblleft{}To peep; to peek\textquotedblright{}.) Say it once more.
\end{tcolorbox}
\section[\texorpdfstring{\ja{まずしい} (mazushii)}{まずしい (mazushii)}]{\ja{まずしい} (mazushii) — needy; poor}

\label{lesson:JA-C304-mazushii}



\begin{quote}
Before the new one: say the Japanese for a trap, then the Japanese for to peep.
\end{quote}

\subsection*{You'll want to know: \ja{まずしい}}
\textbf{\ja{まずしい}} — \emph{mazushii} — \textquotedblleft{}needy; poor\textquotedblright{}.

\begin{quote}
\textbf{\ja{まずしい} \ja{おじいさんと} \ja{おばあさんが} \ja{いました。}} — \emph{mazushii ojiisan to obāsan ga imashita} — \textquotedblleft{}There lived a poor old man and his wife.\textquotedblright{}
\end{quote}

An i-adjective.

\begin{cousinweb}
\textbf{\ja{まずしい}} (\emph{mazushii}) is a native Japanese word (wago), inherited rather than borrowed from Chinese; its kanji was attached later as a kun'yomi (JMdict).
\end{cousinweb}

\begin{grammarlens}[title={What you've built}]
One more word for telling what happened.
\end{grammarlens}

\subsection*{Your turn}
\begin{itemize}
\raggedright
  \item \emph{Say it:} \emph{mazushii}
  \item \emph{Say it:} \emph{mazushii}, once more
  \item \emph{Say it:} \emph{mazushii ojiisan to obāsan ga imashita} — There lived a poor old man and his wife.
  \item \emph{Recall:} say the Japanese for a trap, then the Japanese for to peep, then say \emph{mazushii} again
\end{itemize}

\begin{tcolorbox}[breakable,colback=teal!4,colframe=teal!35!black,title={Before you move on}]
What does \ja{まずしい} mean? (\textquotedblleft{}Needy; poor\textquotedblright{}.) Say it once more.
\end{tcolorbox}
